\chapter{效应修饰}

迄今为止，我们关注的都是整个所关心总体中的平均因果效应。然而，许多因果问题针对的是总体的子集。再考虑那个因果问题：“一个人抬头看天，是否会让其他行人也跟着抬头？”你可能想分别计算“你抬头看天”这一处理在城市居民与外来游客中的平均因果效应，而不是所有行人这个总体中的平均效应。

究竟是在整个总体中还是在其子集中计算平均效应，取决于推断的目的。有些情形下，你并不关心效应在不同个体群体之间的差异。例如，假设你是一位政策制定者，正在考虑是否实施全国性的饮水加氟项目。由于这一公共卫生干预将覆盖总体中的每一个家庭，你主要关心的是整个总体中的平均因果效应，而不是某些特定子集。只有当干预可以针对不同子集实施，或研究结论需要应用于其他总体时，你才会关心因果效应如何在总体的各个子集之间变化。

本章要强调的是：并不存在“处理的因果效应”这种东西。因果效应取决于所研究的特定总体的特征。

\section{处理效应的异质性}

本书开篇计算了心脏移植 $A$ 对死亡 $Y$ 在一个由宙斯大家族 20 名成员构成的总体中的平均因果效应。我们用的是表~1.1 的数据，其各列给出了每个个体（通常观测不到的）反事实结局 $Y^{a=0}$ 与 $Y^{a=1}$ 的取值。考察表~1.1 的数据后，我们得出结论：平均因果效应为零。若所有人都接受心脏移植，总体中有一半人会死亡，$\Pr[Y^{a=1}=1]=10/20=0.5$；若无人接受心脏移植，总体中也有一半人会死亡，$\Pr[Y^{a=0}=1]=10/20=0.5$。因果风险比 $\Pr[Y^{a=1}=1]/\Pr[Y^{a=0}=1]$ 为 $0.5/0.5=1$，因果风险差 $\Pr[Y^{a=1}=1]-\Pr[Y^{a=0}=1]$ 为 $0.5-0.5=0$。

\begin{table}[htbp]
\centering
\caption{宙斯家族成员的反事实结局，按性别 $V$ 分列}
\label{tab:4.1}
\small
\renewcommand{\arraystretch}{1.0}
\begin{tabular}{lccc@{\hspace{1.8em}}lccc}
\toprule
\multicolumn{4}{c}{女性（$V=1$）} & \multicolumn{4}{c}{男性（$V=0$）}\\
\cmidrule(lr){1-4}\cmidrule(lr){5-8}
 & $V$ & $Y^{a=0}$ & $Y^{a=1}$ & & $V$ & $Y^{a=0}$ & $Y^{a=1}$\\
\midrule
Rheia      & 1 & 0 & 1 & Kronos     & 0 & 1 & 0\\
Demeter    & 1 & 0 & 0 & Hades      & 0 & 0 & 0\\
Hestia     & 1 & 0 & 0 & Poseidon   & 0 & 1 & 0\\
Hera       & 1 & 0 & 0 & Zeus       & 0 & 0 & 1\\
Artemis    & 1 & 1 & 1 & Apollo     & 0 & 1 & 0\\
Leto       & 1 & 0 & 1 & Ares       & 0 & 1 & 1\\
Athena     & 1 & 1 & 1 & Hephaestus & 0 & 0 & 1\\
Aphrodite  & 1 & 0 & 1 & Polyphemus & 0 & 0 & 1\\
Persephone & 1 & 1 & 1 & Hermes     & 0 & 1 & 0\\
Hebe       & 1 & 1 & 0 & Dionysus   & 0 & 1 & 0\\
\bottomrule
\end{tabular}
\end{table}

现在我们考虑两个新的因果问题：$A$ 对 $Y$ 在女性中的平均因果效应是多少？在男性中呢？为回答这两个问题，我们将使用表~4.1，它与表~1.1 信息相同，只是多了一列表示性别的指示变量 $V$：$V=1$ 为女性（本书称“女人”），$V=0$ 为男性（本书称“男人”）。\mnote{本章使用“男人”“女人”的说法，可视为一种轻微的用语不当，因为这些神祇是男神与女神，而不是男人与女人。}为方便起见，我们重排了表格，使女性占据前 10 行，男性占据后 10 行。

先计算女性中的平均因果效应。为此，需要把分析限制在表中 $V=1$ 的前 10 行。在这一总体子集中，处理之下的死亡风险为 $\Pr[Y^{a=1}=1\mid V=1]=6/10=0.6$，未处理之下的死亡风险为 $\Pr[Y^{a=0}=1\mid V=1]=4/10=0.4$。因果风险比为 $0.6/0.4=1.5$，因果风险差为 $0.6-0.4=0.2$。也就是说，平均而言，心脏移植 $A$ 增加了女性死亡 $Y$ 的风险。

接着计算男性中的平均因果效应。为此，需要把分析限制在表中 $V=0$ 的后 10 行。在这一总体子集中，处理之下的死亡风险为 $\Pr[Y^{a=1}=1\mid V=0]=4/10=0.4$，未处理之下的死亡风险为 $\Pr[Y^{a=0}=1\mid V=0]=6/10=0.6$。因果风险比为 $0.4/0.6=2/3$，因果风险差为 $0.4-0.6=-0.2$。也就是说，平均而言，心脏移植 $A$ 降低了男性死亡 $Y$ 的风险。

这个例子表明，总体中的平均因果效应为零，并不意味着总体的某个特定子集中的平均因果效应也为零。在表~4.1 中，“无平均因果效应”这一零假设对整个总体成立，但对男性和女性分别而言都不成立。只是男性与女性的平均因果效应大小相等、方向相反。由于两性各占 $50\%$，在考虑整个总体时，两个效应恰好完全抵消。效应的精确抵消大概很罕见，但由于个体对处理的易感性存在差异，处理之个体因果效应的异质性往往是意料之中的。一个例外是严格因果零假设为真之时：此时不存在效应的异质性，因为对每个个体而言效应都为零，从而总体的任何子集中的平均因果效应也都为零。

现在可以给出效应修饰因子的定义了。当 $A$ 对 $Y$ 的平均因果效应随 $V$ 的取值水平而变化时，我们说 $V$ 是 $A$ 对 $Y$ 之效应的一个修饰因子。由于平均因果效应可以用不同的效应度量（如风险差、风险比）来衡量，效应修饰是否存在取决于所用的效应度量。\mnote{效应修饰因子的结构性分类见 6.6 节。\newline 加法效应修饰：层 $V=1$ 与 $V=0$ 的 $\E[Y^{a=1}-Y^{a=0}]$ 不相等。\newline 乘法效应修饰：两层的 $\E[Y^{a=1}]/\E[Y^{a=0}]$ 不相等。\newline 优势比尺度上的效应修饰不予考虑：优势比很少（如果有的话）是因果推断所关心的参数。}例如，性别 $V$ 在加法尺度上是心脏移植 $A$ 对死亡 $Y$ 之效应的修饰因子，因为因果风险差随 $V$ 的水平而变。性别 $V$ 在乘法尺度上也是心脏移植 $A$ 对死亡 $Y$ 之效应的修饰因子，因为因果风险比随 $V$ 的水平而变。我们只把不受处理 $A$ 影响的变量 $V$ 视为效应修饰因子。

在表~4.1 中，因果风险比在女性（$V=1$）中大于 1，在男性（$V=0$）中小于 1。同样，因果风险差在女性（$V=1$）中大于 0，在男性（$V=0$）中小于 0。也就是说，存在\emph{定性效应修饰}，因为子集 $V=1$ 与 $V=0$ 中的平均因果效应方向相反。存在定性效应修饰时，加法效应修饰蕴含乘法效应修饰，反之亦然。然而，不存在定性效应修饰时，可能出现一个尺度（如乘法）上有效应修饰、而另一个尺度（如加法）上没有的情形。为说明这一点，设想在第二项研究中，我们算得边栏中列出的那些量。\mnote{乘法、但非加法的 $V$ 所致效应修饰：\newline $V=1$：\newline $\Pr[Y^{a=0}{=}1]=0.8$，\newline $\Pr[Y^{a=1}{=}1]=0.9$；\newline $V=0$：\newline $\Pr[Y^{a=0}{=}1]=0.1$，\newline $\Pr[Y^{a=1}{=}1]=0.2$。}在这项研究中，$V$ 不引起加法效应修饰，因为 $V=1$ 者的因果风险差等于 $V=0$ 者，即 $0.9-0.8=0.1=0.2-0.1$。但在这项研究中，$V$ 引起了乘法效应修饰，因为 $V=1$ 者的因果风险比不同于 $V=0$ 者，即 $0.9/0.8=1.1\neq 0.2/0.1=2$。由于人们一般不能在不指明所用效应度量（如风险差、风险比）的情况下断言效应修饰存在或不存在，有些作者使用“效应度量修饰”一词而非“效应修饰”，以强调这一概念对效应度量选择的依赖性。

\section{用分层识别效应修饰}

分层分析是识别由已测变量所致效应修饰的自然方法（另见精要~4.1）。要确定 $V$ 是否修饰 $A$ 对 $Y$ 的因果效应，需在变量 $V$ 的每个水平（层）中计算 $A$ 对 $Y$ 的效应。\mnote{分层：$A$ 对 $Y$ 的因果效应在 $V$ 的每一层中分别计算。对二值 $V$，分层后的因果风险差为\newline{\scriptsize$\Pr[Y^{a=1}{=}1\mid V{=}1]$}\newline{\scriptsize${}-\Pr[Y^{a=0}{=}1\mid V{=}1]$,}\newline{\scriptsize$\Pr[Y^{a=1}{=}1\mid V{=}0]$}\newline{\scriptsize${}-\Pr[Y^{a=0}{=}1\mid V{=}0]$}。}上一节中，我们用表~4.1 的数据，计算了移植 $A$ 对死亡 $Y$ 的因果效应在性别 $V$ 的两个层中各自的取值。由于效应在两层之间不同（在加法与乘法尺度上都不同），我们断定 $V$ 对 $A$ 对 $Y$ 的因果效应存在（加法与乘法）效应修饰。

但表~4.1 的数据并不是现实生活中常见的那类数据。现实中你看到的两列，不是每个个体的反事实结局 $Y^{a=1}$ 与 $Y^{a=0}$，而是每个个体的处理水平 $A$ 与观测结局 $Y$。反事实结局不可得，这对用分层来检测效应修饰有何影响？答案取决于研究设计。

先考虑理想的边缘随机化实验。第 2 章已经证明，撇开随机变异不谈，处理的平均因果效应可以用观测数据计算。例如，因果风险差 $\Pr[Y^{a=1}=1]-\Pr[Y^{a=0}=1]$ 等于观测到的关联风险差 $\Pr[Y=1\mid A=1]-\Pr[Y=1\mid A=0]$。同样的推理可以推广到变量 $V$ 的每一层，因为如果处理分配是随机且无条件的，那么总体的每个子集中都应有可交换性。于是，女性中的因果风险差 $\Pr[Y^{a=1}=1\mid V=1]-\Pr[Y^{a=0}=1\mid V=1]$ 等于女性中的关联风险差 $\Pr[Y=1\mid A=1,V=1]-\Pr[Y=1\mid A=0,V=1]$；男性亦然。因此，在无条件随机化的理想实验中，要识别 $V$ 所致的效应修饰，只需做分层分析，即在变量 $V$ 的每一水平上计算关联度量。分层可用于计算总体子集中的平均因果效应，但不能计算个体效应（见精要~2.1 与 3.2）。

现在考虑条件随机化的理想随机化实验。在一个 40 人的总体中，移植 $A$ 以 $0.75$ 的概率随机分配给病情严重者（$L=1$），以 $0.50$ 的概率分配给其他人（$L=0$）。按护照，这 40 人可分为两个国籍：20 名希腊人（$V=1$）与 20 名罗马人（$V=0$）。20 名希腊人的 $L$、$A$ 与死亡 $Y$ 的数据见表~2.2（与表~3.1 相同），20 名罗马人的数据见表~4.2。处理之下的总体风险 $\Pr[Y^{a=1}=1]$ 为 $0.55$，未处理之下的总体风险 $\Pr[Y^{a=0}=1]$ 为 $0.40$。（两个风险都容易用标准化或逆概率加权算出，细节留给读者。）因此，移植 $A$ 对死亡 $Y$ 的平均因果效应在风险差尺度上为 $0.55-0.40=0.15$，在风险比尺度上为 $0.55/0.40=1.375$。在这个总体中，心脏移植增加了死亡风险。

\begin{table}[htbp]
\centering
\caption{罗马人（$V=0$）的观测数据}
\label{tab:4.2}
\small
\renewcommand{\arraystretch}{1.0}
\begin{tabular}{lccc@{\hspace{1.8em}}lccc}
\toprule
\multicolumn{8}{c}{罗马人（$V=0$）}\\
\cmidrule(lr){1-8}
 & $L$ & $A$ & $Y$ & & $L$ & $A$ & $Y$\\
\midrule
Cybele     & 0 & 0 & 0 & Latona     & 1 & 0 & 0\\
Saturn     & 0 & 0 & 1 & Mars       & 1 & 1 & 1\\
Ceres      & 0 & 0 & 0 & Minerva    & 1 & 1 & 1\\
Pluto      & 0 & 0 & 0 & Vulcan     & 1 & 1 & 1\\
Vesta      & 0 & 1 & 0 & Venus      & 1 & 1 & 1\\
Neptune    & 0 & 1 & 0 & Seneca     & 1 & 1 & 1\\
Juno       & 0 & 1 & 1 & Proserpina & 1 & 1 & 1\\
Jupiter    & 0 & 1 & 1 & Mercury    & 1 & 1 & 0\\
Diana      & 1 & 0 & 0 & Juventas   & 1 & 1 & 0\\
Phoebus    & 1 & 0 & 1 & Bacchus    & 1 & 1 & 0\\
\bottomrule
\end{tabular}
\end{table}

如上一章所述，如果这些数据来自一项我们认为条件可交换性 $Y^a\indep A\mid L$ 成立的观察性研究，因果效应的计算也会相同。

下面讨论如何做分层分析，以考察国籍 $V$ 是否修饰 $A$ 对 $Y$ 的效应。目标是计算 $A$ 对 $Y$ 在希腊人中的因果效应 $\Pr[Y^{a=1}=1\mid V=1]-\Pr[Y^{a=0}=1\mid V=1]$，以及在罗马人中的因果效应 $\Pr[Y^{a=1}=1\mid V=0]-\Pr[Y^{a=0}=1\mid V=0]$。如果这两个因果风险差不同，我们就说存在 $V$ 所致的加法效应修饰；若关心乘法效应修饰，对因果风险比作同样处理。

在每个层 $v$ 中计算条件风险 $\Pr[Y^{a=1}=1\mid V=v]$ 与 $\Pr[Y^{a=0}=1\mid V=v]$ 的过程分两阶段：1）按 $V$ 分层；2）按 $L$ 标准化（或等价地，用以 $L$ 为权重的逆概率加权）。我们在第 2 章算过希腊人层（$V=1$）的标准化风险：因果风险差为 0，因果风险比为 1。用同样的过程在罗马人层（$V=0$）中计算，可得风险 $\Pr[Y^{a=1}=1\mid V=0]=0.6$ 与 $\Pr[Y^{a=0}=1\mid V=0]=0.3$。\mnote{当 $V$ 等于达成条件可交换性所需的那些变量 $L$ 时，第 2 步可以省略（见 4.4 节）。}（同样，细节留给读者。）因此，在 $V=0$ 层中，因果风险差为 $0.3$，因果风险比为 2。由于这些效应度量不同于 $V=1$ 层中的相应度量，我们说国籍 $V$ 对移植 $A$ 对死亡 $Y$ 的效应同时存在加法与乘法效应修饰。这一效应修饰不是定性的，因为效应在 $V=0$ 与 $V=1$ 两层中都是有害的或为零。

我们已经表明，在研究总体中，国籍 $V$ 修饰了心脏移植 $A$ 对死亡风险 $Y$ 的效应。但我们并未对这类效应修饰所涉及的因果机制作任何断言。事实上，国籍完全可能只是真正造成效应修饰的那个因果因素的标记。例如，假设希腊的心脏手术质量优于罗马，那么人们就会发现国籍所致的效应修饰。一项改善罗马心脏手术质量的干预，可能消除按护照定义的国籍所致的因果效应修饰。当我们想强调这一区别时，会把国籍称为\emph{替代效应修饰因子}，把医疗质量称为\emph{因果效应修饰因子}。\mnote{见 6.6 节用因果图表示的替代效应修饰因子与因果效应修饰因子。}

因此，我们使用“$V$ 所致的效应修饰”这一说法，并不必然意味着 $V$ 在效应修饰中起因果作用。为避免可能的混淆，有些作者宁愿使用更中性的说法“效应在 $V$ 各层之间的异质性”，而不用“$V$ 所致的效应修饰”。下一章将引入“交互作用”——一个与效应修饰相关、但确实赋予所涉变量以因果作用的概念。

\begin{finepoint}[处理组中的效应]
本章关心总体子集中的平均因果效应。一个特别的子集是受处理者（$A=1$）。若 $\Pr[Y^{a=1}=1\mid A=1]\neq\Pr[Y^{a=0}=1\mid A=1]$，或者由一致性，若 $\Pr[Y=1\mid A=1]\neq\Pr[Y^{a=0}=1\mid A=1]$，则处理组中的平均因果效应不为零。也就是说，若受处理者的观测风险不等于这些受处理者若未受处理时的反事实风险，则处理组中存在因果效应。处理组中的因果风险差为 $\Pr[Y=1\mid A=1]-\Pr[Y^{a=0}=1\mid A=1]$。处理组中的因果风险比又称\emph{标准化发病比}（SMR），为 $\Pr[Y=1\mid A=1]/\Pr[Y^{a=0}=1\mid A=1]$。未处理组中的因果风险差与风险比，可把 $A=1$ 换成 $A=0$ 类似地定义。图~4.1 显示了计算处理组效应与未处理组效应时所比较的组。

若个体因果效应的分布在受处理者与未受处理者之间不同，则处理组中的平均效应将不同于总体中的平均效应。也就是说，计算处理组效应时，处理组 $A=1$ 被当作标记，代表那些真正造成处理组与未处理组之间效应修饰的因素。然而，即便人们可以说，即使处理前变量 $V$ 只是因果效应修饰因子的替代（如国籍），也存在 $V$ 所致的效应修饰，人们也不会说存在处理 $A$ 对效应 $A$ 的修饰，因为那听起来令人困惑。效应修饰是由在处理组之间分布不同的未识别变量所致的。

关于真实效应修饰因子与替代效应修饰因子的图示，见 6.6 节。本书的主体关注总体中的因果效应，因为处理组或未处理组中的因果效应无法直接推广到时变处理（见第三部分）。
\end{finepoint}

\begin{figure}[htbp]
\centering
\begin{tikzpicture}[causal,scale=0.95,every node/.style={transform shape}]
  % 总体
  \begin{scope}
    \clip (0,0.95)--(-0.78,0)--(0,-0.95)--(0.78,0)--cycle;
    \fill[white] (-0.9,-1.0) rectangle (0,1.0);
    \fill[black!30] (0,-1.0) rectangle (0.9,1.0);
  \end{scope}
  \draw[thick] (0,0.95)--(-0.78,0)--(0,-0.95)--(0.78,0)--cycle;
  \draw[thick] (0,0.95)--(0,-0.95);
  \node[font=\small,above=1mm] at (0,0.95) {所关心的总体};
  \node[font=\small,anchor=east] at (-0.95,0) {受治疗};
  \node[font=\small,anchor=west] at (0.95,0) {未受治疗};
  % 箭头
  \draw[->,semithick] (-0.3,-0.88) -- (-2.7,-1.9);
  \draw[->,semithick] (0.3,-0.88) -- (2.7,-1.9);
  \node[font=\small] at (-3.6,-2.15) {处理组的效应};
  \node[font=\small] at (3.6,-2.15) {未处理组的效应};
  % 处理组：白色左半 + 灰色右半
  \fill[white] (-3.95,-2.95)--(-4.6,-3.6)--(-3.95,-4.25)--cycle;
  \draw[thick] (-3.95,-2.95)--(-4.6,-3.6)--(-3.95,-4.25)--cycle;
  \draw[thick] (-3.95,-2.95)--(-3.95,-4.25);
  \fill[black!30] (-3.25,-2.95)--(-3.25,-4.25)--(-2.6,-3.6)--cycle;
  \draw[thick] (-3.25,-2.95)--(-3.25,-4.25)--(-2.6,-3.6)--cycle;
  \node[font=\footnotesize] at (-3.6,-3.6) {vs.};
  % 未处理组：白色左半 + 灰色右半
  \fill[white] (3.25,-2.95)--(2.6,-3.6)--(3.25,-4.25)--cycle;
  \draw[thick] (3.25,-2.95)--(2.6,-3.6)--(3.25,-4.25)--cycle;
  \draw[thick] (3.25,-2.95)--(3.25,-4.25);
  \fill[black!30] (3.95,-2.95)--(3.95,-4.25)--(4.6,-3.6)--cycle;
  \draw[thick] (3.95,-2.95)--(3.95,-4.25)--(4.6,-3.6)--cycle;
  \node[font=\footnotesize] at (3.6,-3.6) {vs.};
  % 标注
  \node[font=\scriptsize,anchor=east] at (-3.85,-4.62) {$\E[Y\mid A=1]$};
  \node[font=\scriptsize,anchor=west] at (-3.3,-4.62) {$\E[Y^{a=0}\mid A=1]$};
  \node[font=\scriptsize,anchor=east] at (3.3,-4.62) {$\E[Y^{a=1}\mid A=0]$};
  \node[font=\scriptsize,anchor=west] at (3.85,-4.62) {$\E[Y\mid A=0]$};
\end{tikzpicture}
\caption{计算处理组效应与未处理组效应时所比较的组。白色与灰色分别代表受治疗者与未受处理者；每一对中，一侧为观测到的结局，另一侧为相应的反事实结局。}
\label{fig:4.1}
\end{figure}

\section{为什么要关心效应修饰}

研究者之所以有兴趣识别效应修饰，并且在随机化实验中也应收集处理前描述变量 $V$ 的数据，有几个相互关联的理由。

第一，如果某因素 $V$ 修饰处理 $A$ 对结局 $Y$ 的效应，那么 $V$ 的患病率不同的总体之间，平均因果效应就会不同。例如，表~4.1 总体中的平均因果效应在女性中有害、在男性中有益，即存在定性效应修饰。由于两性各占 $50\%$，且性别特异的效应一害一益、大小相等而符号相反，整个总体中的平均因果效应为零。然而，假如我们在一个女性比例更高的总体（如应届大学毕业生）中做研究，整个总体的平均因果效应就会是有害的。在存在非定性效应修饰时，平均因果效应的量值（而非方向）可能随总体而变。非定性效应修饰的例子有石棉暴露的效应（在吸烟者与非吸烟者中不同）与全民医疗保障的效应（在低收入家庭与高收入家庭中不同）。

这就是说，一个总体中的平均因果效应取决于该总体中个体因果效应的分布。一般并不存在“处理 $A$ 对结局 $Y$ 的平均因果效应（句号）”这种东西，而只有“在具有特定效应修饰因子组合的总体中，处理 $A$ 对结局 $Y$ 的平均因果效应”。

把一个总体中算得的因果效应外推到第二个总体，称为因果推断在总体之间的\emph{可移植性}（见精要~4.2）。\mnote{有人把缺乏可移植性称为缺乏外部效度。}在我们的例子中，心脏移植 $A$ 对死亡风险 $Y$ 的因果效应在男性与女性之间、罗马人与希腊人之间都不同。因此，这一总体中的平均因果效应未必能移植到效应修饰因子（如性别与国籍）分布不同的其他总体。

由效应修饰因子所定义的各层中的条件因果效应，可能比整个总体中的因果效应更具可移植性；但并不能保证一个总体中的条件效应度量等于另一个总体中的条件效应度量。这是因为可能存在其他未测量或未知的因果效应修饰因子，其条件分布在两个总体之间不同（或因为精要~4.2 中所述的其他原因）。\mnote{Hernán and VanderWeele (2011)、Pearl and Bareinboim (2014)、Dahabreh and Hernán (2019) 等讨论了效应修饰与跨总体推断移植的关系。}这些未测量的效应修饰因子并不是达成可交换性所需的变量，而只是结局的风险因素。因此，效应在总体之间的可移植性，比在单一总体中识别因果效应更难：人们不仅需要对达成可交换性所需的一切因素分层（对这类因素你或许有信息，比如通过访谈那些决定如何分配处理的人），还需要对结局的未测量原因分层，而关于后者的信息要少得多。

因此，因果效应的可移植性是一个不可验证的假设，严重依赖领域知识。例如，多数专家会同意，在尼日尔把家庭年收入增加 \$100 所带来的健康效应（无论加法还是乘法尺度）无法移植到荷兰；但多数专家也会同意，在欧洲人中使用降胆固醇药物的健康效应可以移植到加拿大人。\mnote{一个可移植性也许不成问题的情境：Smith and Pell (2003) 未能找出任何能明显修饰降落伞使用对“重力挑战”（如从高空飞机上跳下）所致死亡之效应的因素。他们认为，在特定人群中开展降落伞使用的随机化试验，不会损害研究结论向其他人群的移植。}

\begin{techpoint}[计算处理组中的效应]
我们在条件可交换性 $Y^a\indep A\mid L$ 下（对 $a=0$ 与 $a=1$ 均成立）计算了总体中的平均因果效应。计算处理组中的平均因果效应，只需要部分可交换性 $Y^{a=0}\indep A\mid L$。换言之，未受处理者若受处理时的风险是否等于实际受处理者的风险，与此无关。未处理组中的平均因果效应则在部分可交换性条件 $Y^{a=1}\indep A\mid L$ 下计算。

下面说明在上述部分可交换性假设下，如何计算形如 $\E[Y^a\mid A=a']$ 的反事实均值。我们用标准化与逆概率加权两种方式来做：
\begin{itemize}
\item 标准化：$\E[Y^a\mid A=a']$ 等于 $\sum_l \E[Y\mid A=a,L=l]\Pr[L=l\mid A=a']$。关于标准化风险比的讨论，见 Miettinen (1972) 与 Greenland and Rothman (2008)。
\item 逆概率加权：$\E[Y^a\mid A=a']$ 等于逆概率加权均值
\[
\frac{\E\!\left[\dfrac{I(A=a)Y}{f(A\mid L)}\Pr[A=a'\mid L]\right]}{\E\!\left[\dfrac{I(A=a)}{f(A\mid L)}\Pr[A=a'\mid L]\right]},
\]
权重为 $\dfrac{\Pr[A=a'\mid L]}{f(A\mid L)}$。对二值 $A$，这一等式由 Sato and Matsuyama (2003) 导出。进一步细节见 Hernán and Robins (2006a)。
\end{itemize}
\end{techpoint}

\begin{finepoint}[可移植性]
在一个总体中估计的效应，常常是为了在另一个总体——目标总体——中作决策。我们能否把效应从研究总体“移植”到目标总体？答案取决于两个总体的特征。具体而言，如果以下特征在两个总体之间相似，因果效应在总体之间的可移植性或许可以得到辩护：
\begin{itemize}
\item 效应修饰：处理的因果效应可能因个体对结局的易感性不同而不同。例如，若女性比男性对处理的效应更易感，我们就说性别是效应修饰因子。一个总体中效应修饰因子的分布，一般会影响该总体中处理之因果效应的量值。若效应修饰因子的分布在两个总体之间不同，则处理的因果效应量值也会不同。
\item 处理的定义不同：处理的因果效应取决于处理在总体中如何定义。若研究总体与目标总体对处理的定义不同，则处理的因果效应量值也会不同。
\item 干扰：正文关注无干扰的情形（精要~1.1）。当一个个体接受处理会影响总体中其他人的结局时，就存在干扰。例如，一个社会活动频繁的人可能说服朋友跟他一起锻炼，于是对该个体身体活动的干预可能比对社交孤立的个体的干预更有效。因此，两个总体之间接触模式的不同，将转化为量值不同的因果效应。
\end{itemize}
越来越多的文献考虑在存在效应修饰时，用研究总体的数据估计目标总体因果效应的可移植性方法（如 Westreich et al. 2017；Rudolph and van der Laan 2017；Dahabreh et al. 2020b）。

把因果推断从一个总体移植到另一个总体，有时可以通过把注意力限定在由效应修饰因子所定义的各层中的平均因果效应，或用研究总体中的层特异效应重构目标总体中的平均因果效应，来加以改善。例如，我们总体中的四个层特异效应度量（罗马女性、希腊女性、罗马男性、希腊男性）可以加权平均，以重构另一个性别与国籍构成不同的总体中的平均因果效应。赋予每个层特异度量的权重，是第二个总体中该层个体所占的比例。然而，由于两个总体之间在未测量效应修饰因子的分布、干扰模式与处理定义上可能存在差异，无法保证这样重构的效应会与目标总体中的真实效应一致。
\end{finepoint}

第二，评估效应修饰是否存在，有助于识别最可能从干预中获益的人群。在表~4.1 的例子中，处理 $A$ 对结局 $Y$ 的平均因果效应为零。然而，处理 $A$ 对男性（$V=0$）有益，对女性（$V=1$）有害。例如，如果医生知道性别存在定性效应修饰，那么在缺乏其他信息的情况下，他们只会治疗碰巧是男性的下一位病人。

当情形如我们的第二个例子那样，存在乘法效应修饰而没有加法效应修饰时，情况稍复杂一些。这里，处理使 $V=0$ 者的结局风险降低 $10\%$，也使 $V=1$ 者的结局风险降低 $10\%$，即 $V$ 不引起加法效应修饰，因为因果风险差在 $V$ 的所有水平上都是 $0.1$。因此，尽管存在乘法效应修饰，一项治疗所有病人的干预在两个 $V$ 层中降低风险的效果相同。事实上，只要 $V$ 的至少一个层中存在非零因果效应，且反事实风险 $\Pr[Y^{a=0}=1\mid V=v]$ 随 $v$ 变化，那么在加法或乘法尺度上必有一处存在效应修饰。

加法（而非乘法）效应修饰才是识别最可能从干预中获益的人群的合适尺度。\mnote{若干作者（如 Blot and Day, 1979；Rothman et al., 1980；Saracci, 1980）认为，加法效应修饰才是公共卫生所关心的那一种。}在不存在加法效应修饰时，得知存在乘法效应修饰对决策恐怕帮助不大。

在我们的第二个例子中，乘法效应修饰的出现是意料之中的，因为 $V=1$ 层中未处理之下的风险等于 $0.8$。于是 $V=1$ 层中可能的最大因果风险比为 $1/0.8=1.25$，它必然不同于 $V=0$ 层中的因果风险比 2。在这类情形中，乘法效应修饰源于未处理之下风险 $\Pr[Y^{a=0}=1\mid V=v]$ 在 $V$ 各水平之间的差异。因此，作为一般规则，报告 $V$ 每一水平 $v$ 上的（绝对）反事实风险 $\Pr[Y^{a=1}=1\mid V=v]$ 与 $\Pr[Y^{a=0}=1\mid V=v]$，比只报告它们的比值或差值更有信息量。

最后，识别效应修饰可能有助于理解导致结局的生物学、社会或其他机制。例如，未受割礼的男性比受割礼者 HIV 感染风险更高，这可能为理解该疾病提供新线索。识别效应修饰也可能成为刻画两种处理之间交互作用的第一步。“效应修饰”与“交互作用”这两个词在科学文献中有时被当作同义词使用。本章关注的是“效应修饰”。下一章把“交互作用”描述为一个因果概念，它与效应修饰相关，但不同于效应修饰。

\section{分层作为一种调整形式}

直到本章之前，我们唯一的目标是计算整个总体中的平均因果效应。在没有边缘随机化时，实现这一目标需要对变量 $L$ 作调整，以保证受处理者与未受处理者的条件可交换性。例如，第 2 章中我们判定心脏移植 $A$ 对死亡 $Y$ 的平均因果效应为零，即因果风险比 $\Pr[Y^{a=1}=1]/\Pr[Y^{a=0}=1]=1$。我们用表~2.2 的数据，通过标准化与逆概率加权两种方法对因素 $L$ 作了调整。

本章为分析增添了另一个可能的目标：识别变量 $V$ 所致的效应修饰。为实现这一目标，除对 $L$ 调整外，还需按 $V$ 分层。例如，本章按国籍 $V$ 分层并对 $L$ 调整，判定心脏移植 $A$ 对死亡 $Y$ 的平均因果效应在希腊人与罗马人之间不同。

总之，标准化（或逆概率加权）用于对 $L$ 调整，分层用于识别 $V$ 所致的效应修饰。

但分层并不总是用于识别 $V$ 所致的效应修饰。实践中，分层常被当作标准化（及逆概率加权）的替代方法来对 $L$ 调整。事实上，用分层对 $L$ 调整的做法如此普遍，以致许多研究者把“分层”与“调整”视为同义词。例如，假设你请一位流行病学家对因素 $L$ 调整，以计算心脏移植 $A$ 对死亡 $Y$ 的效应。她多半会立刻把表~2.2 拆成两个子表——一个限于 $L=0$ 的个体，另一个限于 $L=1$ 的个体——并给出各自的效应度量（比如风险比）。也就是说，她会算出风险比 $\Pr[Y=1\mid A=1,L=l]/\Pr[Y=1\mid A=0,L=l]=1$，对 $l=0$ 与 $l=1$ 都成立。在给定 $L$ 的条件可交换性下，这两个层特异的关联风险比可以赋予因果解释：它们分别度量由 $L=0$ 与 $L=1$ 定义的总体子集中的平均因果效应。它们是\emph{条件效应度量}。相比之下，我们在第 2 章算得的风险比 1 是一个\emph{边缘}（无条件）效应度量。\mnote{在给定 $L$ 的条件可交换性下，子集 $L=l$ 中的风险比度量该子集的平均因果效应：若 {\scriptsize$Y^a\indep A\mid L$}，则观测风险等于反事实风险，即\newline{\scriptsize$\Pr[Y{=}1\mid A{=}a,L{=}l]$}\newline{\scriptsize${}=\Pr[Y^a{=}1\mid L{=}l]$}。}在这个特定例子中，三个风险比——两个条件的与一个边缘的——恰好相等，因为不存在 $L$ 所致的效应修饰。分层必然产生多个层特异的效应度量（由变量 $L$ 定义的每一层一个）。每一个都量化总体中一个互不重叠的子集里的平均因果效应，但一般而言，它们没有一个量化整个总体中的平均因果效应。因此，第 2 章描述计算总体中处理之平均因果效应的方法时，我们没有考虑分层，而是着眼于标准化与逆概率加权。

此外，与标准化和逆概率加权不同，通过分层作调整需要在由条件可交换性所需全部变量 $L$ 的各种组合所定义的总体子集中计算效应度量。例如，用分层估计表~2.2 与 4.2 总体中心脏移植的效应时，必须计算 $L=1$ 的罗马人、$L=1$ 的希腊人、$L=0$ 的罗马人、$L=0$ 的希腊人中的效应；但不能只计算 $V=0$ 层中的关联就得到罗马人中的效应，因为国籍 $V$ 本身不足以保证条件可交换性。

也就是说，使用分层迫使人们评估条件可交换性所需的所有变量 $L$ 所致的效应修饰，无论你是否对这种效应修饰感兴趣。相反，先按 $V$ 分层、再用逆概率加权或标准化对 $L$ 调整，则可以如上所述，把可交换性与效应修饰分开处理。

与使用分层相关的其他问题还有：某些效应度量（如优势比）的\emph{不可压缩性}（见精要~4.3）；以及在时变处理的情形下，当需要对受先前处理影响的时变变量 $L$ 调整时，不恰当的调整会导致偏倚（见第三部分）。\mnote{时变处理下，层特异效应度量未必有因果解释（Robins 1986, 1987；见第 20 章）。\newline 分层还要求正值性：只含受处理者或只含未受处理者的层 $L=l$ 中，因果效应无法计算。}

有时研究者只对变量 $L$ 所定义的某些层计算因果效应，即对某些层不计算层特异的效应度量。这种形式的分层称为\emph{限制}。对因果推断而言，分层不过是对总体的若干个穷尽且互斥的子集应用限制，并在每个子集内都具备可交换性。当总体中某些层的正值性不成立时，用限制把因果推断限定在原始总体中正值性成立的那些层（见第 3 章）。

\section{匹配作为另一种调整形式}

匹配是另一种调整方法。匹配的目标是构造总体的一个子集，使变量 $L$ 在受处理者与未受处理者中分布相同。\mnote{匹配的讨论只适用于队列研究。病例对照设计（见第 8 章）中匹配的是病例与对照，而非受处理者与未受处理者；即便匹配因子足以保证条件可交换性，这样做也不能获得匹配总体中两组的无条件可交换性，仍需按匹配因子分层以估计条件（层特异）效应度量。}举例来说，取表~2.2 的心脏移植例子，其中变量 $L$ 足以达成条件可交换性。对每个病情非危重（$A=0,L=0$）的未受处理者，随机选取一名病情非危重的受处理者（$A=1,L=0$）；对每个病情危重（$A=0,L=1$）的未受处理者，随机选取一名病情危重的受处理者（$A=1,L=1$）。我们把每个未受处理者及其对应的受处理者称为一个匹配对，把变量 $L$ 称为匹配因子。设我们构成如下 7 个匹配对：$L=0$ 的 Rheia-Hestia、Kronos-Poseidon、Demeter-Hera、Hades-Zeus，以及 $L=1$ 的 Artemis-Ares、Apollo-Aphrodite、Leto-Hermes。总体中所有未受处理者都被选中，而受处理者只抽取了一部分。在这个由匹配对构成的总体子集中，病情危重者（$L=1$）的比例按设计在受处理者与未受处理者中相同（$3/7$）。

为构造匹配总体，我们用受处理者的一个子集替换了总体中的受处理者，该子集中匹配因子 $L$ 的分布与未受处理者相同。在给定 $L$ 的条件可交换性假设下，这一过程的结果是：匹配总体中受处理者与未受处理者（无条件）可交换。由于匹配总体中二者可交换，其平均结局可以直接比较：受处理者的风险为 $3/7$，未受处理者的风险为 $3/7$，故因果风险比为 1。注意，匹配保证匹配总体中的正值性，因为只含受处理者或只含未受处理者的层被排除在分析之外。

人们常选择个体较少的那一组（我们例子中的未受处理者），用另一组（我们例子中的受处理者）为其寻找匹配。被选中的那一组定义了计算因果效应所在的子总体。上一段我们计算的是未处理组中的效应。若在所有 $L$ 层中受处理者都少于未受处理者，我们通常计算处理组中的效应。此外，匹配不必是一对一的（匹配对），也可以是一对多的（匹配集）。

在许多应用中，$L$ 是若干变量构成的向量。此时，对由 $L$ 中所有变量取值组合所定义的某一层中的每个未受处理者，我们会从同一层中随机选取一名（或数名）受处理者。\mnote{匹配因子越多，个体找到精确匹配的可能性越小；寻找近似匹配的方法见 Stuart (2010) 的导论。}

匹配可用于构造具有任意选定 $L$ 分布的匹配总体，而不限于受处理者或未受处理者中的分布。感兴趣的分布可以通过如上所述的个体匹配实现，也可以通过\emph{频数匹配}实现。后者的一个例子是：随机抽取受处理者，使其中 $70\%$ 的人有 $L=1$，然后对未受处理者重复同样的程序。

由于匹配总体是原研究总体的一个子集，匹配后研究总体中因果效应修饰因子的分布，一般会不同于原本未匹配的研究总体，如下一节所述。

\section{效应修饰与调整方法}

标准化、逆概率加权、分层/限制与匹配是估计平均因果效应的不同方法，但它们估计的是不同类型的因果效应。这四种方法可按所估计效应的类型分为两组：标准化与逆概率加权可用于计算边缘效应或条件效应，分层/限制与匹配则只能用于计算总体中某些子集的条件效应。四种方法都要求可交换性与正值性，但这些条件需要成立的总体子集取决于所关心的因果效应。例如，要计算 $L=l$ 者的条件效应，上述任一方法只要求该子集中的可交换性与正值性；要估计整个总体中的边缘效应，则要求可交换性与正值性在 $L$ 的所有水平上成立。

不存在效应修饰时，用这四种方法算得的效应度量（风险比或风险差）将相等。例如，我们曾断定心脏移植 $A$ 对死亡 $Y$ 的平均因果效应为零：无论在表~2.2 的整个总体中（标准化与逆概率加权），在病情危重 $L=1$ 与病情非危重 $L=0$ 的总体子集中（分层），还是在未受处理者中（匹配），结果都是如此。所有方法给出的因果风险比都等于 1。然而，用这四种方法算得的效应度量一般并不相等。为说明效应可以如何不同，我们用表~4.3 的数据计算心脏移植 $A$ 对高血压 $Z$（1：有，0：无）的效应。我们假定可交换性 $Z^a\indep A\mid L$ 与正值性成立。选用风险比尺度并无特别理由。

\begin{table}[htbp]
\centering
\caption{心脏移植 $A$、病情严重程度 $L$ 与高血压 $Z$ 的观测数据}
\label{tab:4.3}
\small
\renewcommand{\arraystretch}{1.0}
\begin{tabular}{lccc@{\hspace{1.8em}}lccc}
\toprule
 & $L$ & $A$ & $Z$ & & $L$ & $A$ & $Z$\\
\midrule
Rheia      & 0 & 0 & 0 & Leto       & 1 & 0 & 0\\
Kronos     & 0 & 0 & 1 & Ares       & 1 & 1 & 1\\
Demeter    & 0 & 0 & 0 & Athena     & 1 & 1 & 1\\
Hades      & 0 & 0 & 0 & Hephaestus & 1 & 1 & 1\\
Hestia     & 0 & 1 & 0 & Aphrodite  & 1 & 1 & 0\\
Poseidon   & 0 & 1 & 0 & Polyphemus & 1 & 1 & 0\\
Hera       & 0 & 1 & 1 & Persephone & 1 & 1 & 0\\
Zeus       & 0 & 1 & 1 & Hermes     & 1 & 1 & 0\\
Artemis    & 1 & 0 & 1 & Hebe       & 1 & 1 & 0\\
Apollo     & 1 & 0 & 1 & Dionysus   & 1 & 1 & 0\\
\bottomrule
\end{tabular}
\end{table}

标准化与逆概率加权给出整个总体中平均因果效应的风险比 $\Pr[Z^{a=1}=1]/\Pr[Z^{a=0}=1]=0.8$（这些及以下计算留给读者）。分层给出条件因果风险比：在 $L=0$ 层中 $\Pr[Z^{a=1}=1\mid L=0]/\Pr[Z^{a=0}=1\mid L=0]=2.0$，在 $L=1$ 层中 $\Pr[Z^{a=1}=1\mid L=1]/\Pr[Z^{a=0}=1\mid L=1]=0.5$。用上一节选出的匹配对作匹配，给出未处理组中的因果风险比 $\Pr[Z^{a=1}=1\mid A=0]/\Pr[Z=1\mid A=0]=1.0$。

我们算出了四个因果风险比，得到四个不同的数：$0.8$、$2.0$、$0.5$ 与 $1.0$。它们全都正确。撇开随机变异不谈（见技术注记~4.2），差异的解释是定性效应修饰：处理使病情非危重者（$L=0$）的风险加倍（因果风险比 $2.0$），使病情危重者（$L=1$）的风险减半（因果风险比 $0.5$）。总体中的平均因果效应（因果风险比 $0.8$）是有益的，因为危重组未处理之下的反事实风险与非危重组之比 $\Pr[Z^{a=0}=1\mid L=1]/\Pr[Z^{a=0}=1\mid L=0]$，超过了处于非危重组的优势 $\Pr[L=0]/\Pr[L=1]$ 的 2 倍（见技术注记~4.3）。未处理组中的因果效应为零（因果风险比 $1.0$），这反映了未处理者中病情非危重者的比例高于整个总体。这个例子凸显了首要的一点：必须明确效应度量所对应的总体，或总体的某个子集。

上一章主张，定义良好的因果效应是有意义的因果推断的一个前提。本章则主张，刻画良好的目标总体是另一个前提。在比较两种或多种干预、且总体满足某些事先设定的合格标准的实验中，这两个前提自动具备。但在观察性研究中，这些前提不能想当然。相反，开展观察性研究的研究者需要明确界定所关心的因果效应，以及计算该效应所在的总体子集。否则，当用不同方法得到的效应度量不同时，很容易产生误解。

在上面的例子中，一位使用逆概率加权（计算整个总体中的效应）的研究者与另一位使用匹配（计算未处理组中的效应）的研究者，不必就哪种分析方法更优越展开争论。他们效应度量的差异源于各自提出的因果问题不同，而不是源于分析方法的选择。事实上，第二位研究者本可以用逆概率加权计算未处理组或处理组中的效应（见技术注记~4.1）。

最后一点。分层可用于计算总体子集中的平均因果效应，但不能计算个体（个体特异）效应。如前所述，个体因果效应只能在极端假设下识别。见精要~2.1 与 3.2。

\begin{techpoint}[层特异效应度量的合并]
直到第 10 章，我们都通过假定自己处理的是整个总体而非样本来回避统计上的考虑。因此我们谈的是计算因果效应，而不是（一致地）估计它们。然而在实践中，我们很少能计算总体中的因果效应。我们从样本中估计它们，并希望围绕效应估计得到相当窄的置信区间。

处理层特异的效应度量时，减小估计变异性的一个常见做法，是把所有层特异的效应度量合并成一个合并的层特异效应度量。其想法是：若不存在效应度量修饰，则合并后的效应度量将比任何单个层特异效应度量都更精确地估计那个共同的效应度量。合并方法（如 Woolf、Mantel-Haenszel、最大似然）有时计算层特异效应度量的加权平均，权重选取得使合并估计的变异性最小。Greenland and Rothman (2008) 综述了分层分析的一些常用方法。合并的效应度量也可以用回归模型计算，模型中包含所有协变量 $L$ 之间的一切可能乘积项，但不含处理 $A$ 与协变量 $L$ 之间的乘积项，即模型对 $L$ 是饱和的（见第 11 章）。

合并的主要目的是围绕共同的层特异效应度量得到更窄的置信区间，但合并后的效应度量仍然是条件效应度量。在我们的心脏移植例子中，结局 $Z$ 的合并层特异风险比（Mantel-Haenszel 法）为 $0.88$。这一结果只有在层特异风险比 2 与 $0.5$ 确实是同一层特异因果效应的估计时才有意义。例如，假设两层的因果风险比都是 $0.9$，但由于样本量小，我们得到 $0.5$ 与 $2.0$ 的估计。那种情况下合并是恰当的，Mantel-Haenszel 风险比会比任一层的风险比更接近真相。反之，若层特异的因果风险比确实是 $0.5$ 与 $2.0$，那么合并就没有什么意义，Mantel-Haenszel 风险比也不易解释。对处理效应异质的研究做 meta 分析时也会遇到同样的问题（Dahabreh et al. 2020a）。

在实践中，判断各层之间效应度量的异质性究竟源于抽样变异还是源于效应度量修饰，并不总是显而易见。分层越细，随机变异引入的不确定性越大。
\end{techpoint}

\begin{techpoint}[边缘与条件因果风险比之间的关系]
假设我们想知道，在已知每一层 $l$ 的条件风险比 $\Pr[Y^{a=1}=1\mid L=l]/\Pr[Y^{a=0}=1\mid L=l]$ 的取值时，边缘风险比 $\Pr[Y^{a=1}=1]/\Pr[Y^{a=0}=1]$ 在什么条件下会小于 1。为此注意
\begin{gather*}
\frac{\Pr[Y^{a=1}=1]}{\Pr[Y^{a=0}=1]}=
\sum_l \frac{\Pr[Y^{a=1}=1\mid L=l]}{\Pr[Y^{a=0}=1\mid L=l]}\,w(l),\\
w(l)=\frac{\Pr[Y^{a=0}=1\mid L=l]\Pr[L=l]}{\Pr[Y^{a=0}=1]},
\end{gather*}
且 $\sum_l w(l)=1$。把 $w(1)$ 与 $w(0)$ 代入，再作一些代数运算，即可得到不等式 $\Pr[Y^{a=1}=1]/\Pr[Y^{a=0}=1]<1$ 成立的条件。

在我们的数据例子中，$\Pr[Y^{a=1}=1\mid L=l]/\Pr[Y^{a=0}=1\mid L=l]$ 在 $L=1$ 时为 $0.5$，在 $L=0$ 时为 $2.0$。因此，边缘风险比小于 1，当且仅当
\[
\frac{\Pr[Y^{a=0}=1\mid L=1]}{\Pr[Y^{a=0}=1\mid L=0]}>2\,\frac{\Pr[L=0]}{\Pr[L=1]}.
\]
\end{techpoint}

\begin{table}[htbp]
\centering
\caption{20 名基线无抑郁者的观测数据}
\label{tab:4.4}
\small
\renewcommand{\arraystretch}{1.0}
\begin{tabular}{lccc@{\hspace{1.8em}}lccc}
\toprule
 & $V$ & $A$ & $Y$ & & $V$ & $A$ & $Y$\\
\midrule
Rheia      & 1 & 0 & 0 & Kronos     & 0 & 0 & 0\\
Demeter    & 1 & 0 & 0 & Hades      & 0 & 0 & 0\\
Hestia     & 1 & 0 & 0 & Poseidon   & 0 & 0 & 1\\
Hera       & 1 & 0 & 0 & Zeus       & 0 & 0 & 1\\
Artemis    & 1 & 0 & 1 & Apollo     & 0 & 0 & 0\\
Leto       & 1 & 1 & 0 & Ares       & 0 & 1 & 1\\
Athena     & 1 & 1 & 1 & Hephaestus & 0 & 1 & 1\\
Aphrodite  & 1 & 1 & 1 & Polyphemus & 0 & 1 & 1\\
Persephone & 1 & 1 & 0 & Hermes     & 0 & 1 & 0\\
Hebe       & 1 & 1 & 1 & Dionysus   & 0 & 1 & 1\\
\bottomrule
\end{tabular}
\end{table}

\begin{finepoint}[可压缩性与优势比]
在不存在 $V$ 所致乘法效应修饰时，整个总体中的因果风险比 $\Pr[Y^{a=1}=1]/\Pr[Y^{a=0}=1]$ 等于 $V$ 每一层 $v$ 中的条件因果风险比 $\Pr[Y^{a=1}=1\mid V=v]/\Pr[Y^{a=0}=1\mid V=v]$。更一般地，因果风险比是各层特异风险比的加权平均。例如，若 $V=1$ 与 $V=0$ 两层的因果风险比分别为 2 与 3，则总体中的因果风险比将大于 2 而小于 3。总体中因果风险比（以及因果风险差）的取值总被层特异风险比的取值范围所约束，这一点不仅显然，而且是任何效应度量都应具备的良好性质。

现在考虑一种假想的效应度量（不是风险比或风险差），它使得总体效应度量不是各层特异度量的加权平均。也就是说，总体效应度量未必落在层特异效应度量的取值范围内。这样的效应度量会是很古怪的。优势比（双关语）正是这样一种效应度量，下面我们来讨论。

假设表~4.4 的数据是为计算海拔 $A$ 对抑郁 $Y$ 的因果效应而收集的，总体为 20 名基线时无抑郁者。若个体迁往高海拔住所（奥林匹斯山顶），处理 $A$ 为 1，否则为 0；若个体随后出现抑郁，结局 $Y$ 为 1，否则为 0；若个体为女性，$V$ 为 1，男性为 0。迁移的决定是随机的，即那些更易患抑郁的人与其他人一样可能迁移；实际上 $Y^a\indep A$。因此，风险比 $\Pr[Y=1\mid A=1]/\Pr[Y=1\mid A=0]=2.3$ 就是总体中的因果风险比，而优势比
\[
\frac{\Pr[Y=1\mid A=1]/\Pr[Y=0\mid A=1]}{\Pr[Y=1\mid A=0]/\Pr[Y=0\mid A=0]}=5.4
\]
就是总体中的因果优势比
\[
\frac{\Pr[Y^{a=1}=1]/\Pr[Y^{a=1}=0]}{\Pr[Y^{a=0}=1]/\Pr[Y^{a=0}=0]}.
\]
风险比与优势比在不同的尺度上度量同一因果效应。

现在计算性别特异的因果效应，分别在风险比与优势比尺度上。条件因果风险比 $\Pr[Y=1\mid V=v,A=1]/\Pr[Y=1\mid V=v,A=0]$ 在男性（$V=0$）中为 2，在女性（$V=1$）中为 3。条件因果优势比
\[
\frac{\Pr[Y=1\mid V=v,A=1]/\Pr[Y=0\mid V=v,A=1]}{\Pr[Y=1\mid V=v,A=0]/\Pr[Y=0\mid V=v,A=0]}
\]
在男性（$V=0$）中为 6，在女性（$V=1$）中为 6。总体中的因果风险比 $2.3$ 介于性别特异的因果风险比 2 与 3 之间。相反，总体中的因果优势比 $5.4$ 比两个性别特异的优势比 6 都小（即更接近零值）。以优势比尺度衡量，因果效应在总体的每一半中都大于整个总体。当 $V$ 是 $Y$ 的独立风险因素，且如我们的随机化实验那样 $A$ 独立于 $V$ 时（Miettinen and Cook, 1981），总体因果优势比可能比非零的层特异因果优势比更接近零值。

当总体效应度量可以表示为层特异度量的加权平均时，我们说该效应度量是\emph{可压缩的}。在随访研究中，风险比与风险差是可压缩的效应度量，而优势比——或很少使用的优势差——不是（Greenland 1987）。优势比的不可压缩性是 Jensen 不等式的一个特例（Samuels 1981），它可能导致上述那类反直觉的结果。优势比在严格因果零假设下是可压缩的——此时条件与无条件效应度量都等于零值——并且当结局在随访研究的每一层中都罕见（比如 $<10\%$）时，它近似可压缩，且近似等于风险比。

优势比不可压缩性的一个重要推论是：在逻辑上不可能把“缺乏可交换性”与“条件优势比相对于无条件优势比发生变化”等同起来。在我们的例子中，尽管受处理者与未受处理者是可交换的，优势比却变化了约 $10\%$（$1-6/5.4$）。Greenland、Robins 与 Pearl (1999) 综述了不可压缩性与缺乏可交换性之间的关系。
\end{finepoint}
